\documentclass[11pt]{article}
\usepackage[T1]{fontenc}
\usepackage[utf8]{inputenc}
\usepackage{amsmath,amssymb,amsthm,mathtools}
\usepackage{geometry,graphicx,booktabs,longtable,tabularx,array}
\usepackage{enumitem,needspace,xcolor,tikz,tcolorbox}
\usetikzlibrary{arrows.meta,positioning,calc,fit,backgrounds,decorations.pathreplacing,matrix,patterns}
\tcbuselibrary{breakable}
\usepackage{caption,hyperref,bookmark,cleveref,microtype}
\usepackage{komorebi-publication}
\renewcommand{\KomorebiTitle}{Komorebi Space}
\renewcommand{\KomorebiSubtitle}{Coherent realizations, obstructions, and descent}
\renewcommand{\KomorebiSeries}{Foundations of the Realization Calculus}
\renewcommand{\KomorebiKind}{Monograph}
\renewcommand{\KomorebiEdition}{Expanded monograph edition 1.4.0 / 11 October 2026}
\renewcommand{\KomorebiShortTitle}{Komorebi Space}
\renewcommand{\KomorebiProfile}{monograph}
\KomorebiSetup
\setcounter{tocdepth}{1}
\setcounter{secnumdepth}{2}
\setlist{nosep,leftmargin=1.7em,topsep=4pt}
\hypersetup{pdfkeywords={Komorebi Space,realization calculus,descent,obstruction,Lean,holonomy,closed range}}
\newtheorem{theorem}{Theorem}[section]
\newtheorem{proposition}[theorem]{Proposition}
\newtheorem{lemma}[theorem]{Lemma}
\newtheorem{corollary}[theorem]{Corollary}
\theoremstyle{definition}
\newtheorem{definition}[theorem]{Definition}
\newtheorem{example}[theorem]{Example}
\theoremstyle{remark}
\newtheorem{remark}[theorem]{Remark}
\numberwithin{equation}{section}
\newcommand{\R}{\mathbb R}
\newcommand{\Z}{\mathbb Z}
\newcommand{\Q}{\mathbb Q}
\newcommand{\N}{\mathbb N}
\newcommand{\K}{\mathbb K}
\newcommand{\V}{\mathcal V}
\newcommand{\W}{\mathcal W}
\newcommand{\id}{\operatorname{id}}
\newcommand{\im}{\operatorname{im}}
\newcommand{\ran}{\operatorname{ran}}
\newcommand{\Hom}{\operatorname{Hom}}
\newcommand{\Ob}{\operatorname{Ob}}
\newcommand{\Ext}{\operatorname{Ext}}
\newcommand{\Match}{\operatorname{Match}}
\newcommand{\Exact}{\operatorname{Exact}}
\newcommand{\dist}{\operatorname{dist}}
\newcommand{\Aut}{\operatorname{Aut}}
\newcommand{\Stab}{\operatorname{Stab}}
\newcommand{\Set}{\mathbf{Set}}
\newcommand{\Cat}{\mathbf{Cat}}
\newcommand{\norm}[1]{\left\lVert#1\right\rVert}
\newcommand{\abs}[1]{\left\lvert#1\right\rvert}
\newcommand{\source}[1]{\par\phantomsection\label{src:#1}{\footnotesize\sffamily\color{KomorebiMuted}Source concordance: \hyperlink{claim:#1}{#1}.}\par\smallskip}
\newcommand{\code}[1]{\texttt{\detokenize{#1}}}
\newcommand{\chap}[1]{\par\Needspace{10\baselineskip}\section{#1}}
\newcommand{\partnote}[2]{\clearpage\part{#1}\noindent\begin{minipage}{.92\linewidth}\large #2\end{minipage}\par\vspace{12pt}}
\newtcolorbox{insight}{breakable,colback=KomorebiSoft!55,colframe=KomorebiPine!50,boxrule=.45pt,arc=1pt,left=9pt,right=9pt,top=7pt,bottom=7pt}
\newtcolorbox{scopebox}{breakable,colback=KomorebiWarm,colframe=KomorebiOchre!55,boxrule=.4pt,arc=1pt,left=9pt,right=9pt,top=7pt,bottom=7pt}
\newcommand{\plate}[3]{\begin{figure}[!htbp]\centering\resizebox{.96\linewidth}{!}{\input{figures/#1}}\caption{#2}\label{#3}\end{figure}}
\tikzset{box/.style={draw=KomorebiPine!70,rounded corners=2pt,fill=KomorebiSoft!50,align=center,inner sep=8pt,font=\small},arr/.style={-{Latex[length=2.2mm]},draw=KomorebiPine,thick},thinarr/.style={-{Latex[length=1.8mm]},draw=KomorebiMuted},lab/.style={font=\small,align=center},dot/.style={circle,fill=KomorebiPine,inner sep=2.3pt}}
\renewcommand{\arraystretch}{1.18}
\raggedbottom
\widowpenalty=7000\clubpenalty=7000
\setlength{\emergencystretch}{3em}

% Compact only the contents entries, without reducing body type size.
\usepackage{etoolbox}
\makeatletter
\patchcmd{\l@section}{1.0em}{0.6em}{}{}
\makeatother

% Generated by checks/edition_stats.py from the current executed reports.
% Generated from current executed reports. Do not hand edit.
\newcommand{\EditionClaims}{80}
\newcommand{\EditionReferences}{211}
\newcommand{\EditionWrittenOnly}{0}
\newcommand{\EditionOriginalChecks}{14,224}
\newcommand{\EditionRecoveredChecks}{2,767}
\newcommand{\EditionRevisionChecks}{17,213}
\newcommand{\EditionExpansionChecks}{96,345}
\newcommand{\EditionDeepeningChecks}{13,817}
\newcommand{\EditionTotalChecks}{144,366}
\newcommand{\EditionFormalModules}{138}
\newcommand{\EditionAuditedDeclarations}{4,745}
\newcommand{\EditionNewTheorems}{54}
\newcommand{\EditionFormalTheorems}{1,507}
\newcommand{\EditionNewConstructions}{17}
\newcommand{\EditionReplays}{1}
\newcommand{\EditionNegativeControls}{7}


% Keep readable half-page figures with their surrounding exposition.
\renewcommand{\topfraction}{.88}
\renewcommand{\bottomfraction}{.80}
\renewcommand{\textfraction}{.10}
\renewcommand{\floatpagefraction}{.78}
\setcounter{topnumber}{2}
\setcounter{bottomnumber}{2}
\setcounter{totalnumber}{3}

\begin{document}
\KomorebiCover
\pagenumbering{roman}
\section*{About this edition}
This monograph develops Komorebi Space as a theory of coherent realizations relative to a specified system of mathematical views. Its central constructions are ordinary limits, relative extension spaces, obstruction quotients, cocycle groupoids, an explicit stack on a composition site, and quantitative repair operators. The point of the synthesis is to keep these constructions connected without erasing the differences between their hypotheses.

The exposition follows the mathematics rather than the order in which the source library grew. Definitions precede existence questions. Existence is separated from uniqueness. Algebraic feasibility is separated from numerical stability. Changing a representation is separated from changing the data that the representation is required to preserve.

Publication version \textbf{1.4.0} develops ten further principal exposition items from the unchanged formal baseline. A new part explains exactification with variable rates, the difference between an algebraic obstruction and its numerical floor, and a complete infinite-dimensional model where exact finite repairs escape the intended Hilbert space. Additional chapters derive a directed geometry from a separately declared transport cost and work through a coherent phase that cannot be removed by scalar changes. The earlier finite-observation, repair-comparison, and loop developments remain intact.

The canonical Lean baseline is version 1.2.0. It contains \EditionFormalModules\ production modules and one root. Its accepted clean replay and imported axiom audit cover \EditionAuditedDeclarations\ project declarations, including generated declarations. This edition retains those records and checks that every formal file is unchanged. It does \emph{not} claim a new Lean compilation. Fresh checks in this edition cover source correspondence, exact finite models, figure generation, and the publication build. Historical archives and compiler records retain their original dates and scopes.

Lean's compiler, kernel, and pinned external dependency cache remain trusted. This is not an independent rebuild of mathlib or a proof that a Python program implements every formal construction. Source correspondence identifies the actual declarations and hypotheses. It does not turn explanatory prose into a kernel-checked language.

All mathematical figures are reproducible vector diagrams or Python plots of explicit formulas. Axes, support patterns, and worked values are checked against their stated mathematical models. The source package contains the manuscript, figure sources, exact checks, evidence archives, and the progress ledger. The separate visual guide offers a shorter route through the examples. The Lean guide explains the correspondence and its verification boundaries.

\begin{scopebox}
\textbf{Reading scope.} The word \emph{space} does not assert one universal topological object containing all of mathematics. The strict construction is a set of coherent families. Symmetry is retained in a groupoid, and effective descent is proved for a specified site. Analytic estimates apply only after a normed or Hilbert model has been declared. These are related layers, not interchangeable definitions.
\end{scopebox}

\clearpage
\section*{Orientation}
A description at one scale often omits information visible at another. A formula in one representation may fail to commute with a change of representation. Local pieces may each be valid and still fail to assemble. Even when assembly is possible, there may be many completions, and small discrepancies may or may not imply proximity to an exact answer.

Komorebi Space treats these questions in a common order. First specify the views and the maps between them. Then define the coherent realizations. Next identify the information that must remain fixed. Finally determine the obstruction to extension, the freedom in a successful extension, and the conditions under which repair is stable.

The resulting program has three central distinctions:
\[
\begin{gathered}
\text{a collection of appearances}\quad\ne\quad\text{a coherent realization},\\
\text{existence of an exact answer}\quad\ne\quad\text{a stable repair estimate},\\
\text{one isomorphism class}\quad\ne\quad\text{absence of symmetries}.
\end{gathered}
\]
The monograph makes these distinctions into explicit constructions and theorems.

\subsection*{The route through the book}
Part I constructs coherent sections and follows their obstructions as more data are prescribed. Part II develops resolution towers, finite observations, split repair coordinates, and multiscale reconstruction. Part III retains symmetries, proves effective descent, and distinguishes coherent phases from removable ones. Part IV develops analytic stability and an independent geometry of transformation costs. Part V follows approximate answers to their limits and exhibits the exact boundary between a coherent point and an intended Hilbert realization. Part VI works through two finite loop models and returns to mathematical semantics. Appendix~\ref{app:workshop} contains nine solved problems, and Appendix~\ref{app:hypotheses} gives a hypothesis atlas.

A short technical route is Sections~\ref{sec:views}, \ref{sec:obstruction}, \ref{sec:stack}, \ref{sec:closed}, and~\ref{sec:motion}. A reader interested in the role of packets and residuals should include Sections~\ref{sec:split} and~\ref{sec:cascade}. The graph models in Section~\ref{sec:graph} explain how a categorical realization problem acquires an actual numerical condition number. Sections~\ref{sec:observation} and~\ref{sec:repair-transport} give the representation-theoretic links. Sections~\ref{sec:signal-loop}, \ref{sec:twisted-loop}, and~\ref{sec:escape} can also be read first as concrete motivations. The last of these connects finite realization, closure, and failure of protected infinite existence in one calculation.

\subsection*{What is being contributed}
The limit universal property, gauge action groupoids, standard stack definitions, orthogonal projection, and bounded inverse principles belong to established mathematics. They are used here as such. The contribution of this presentation is their explicit integration with the project's actual view systems, protected correction spaces, resolution differentials, graph transports, and checked examples. Historical novelty of every component is not inferred from formal verification. Conversely, the fact that the ingredients are classical does not remove the need to prove the comparison maps and to retain the right hypotheses.

\clearpage\begingroup\setlength{\parskip}{0pt}\tableofcontents\endgroup
\clearpage\pagenumbering{arabic}
\input{chapters/01_views}
\input{chapters/02_obstructions}
\input{chapters/02a_prescription}
\input{chapters/03_towers}
\input{chapters/04_symmetry}
\input{chapters/04a_boundary_lab}
\clearpage
\chap{Coherence need not mean removable phase}\label{sec:phase}
The preceding descent theory starts with transition data that already satisfy their composition laws. It is useful to see, in a completely finite calculation, how those laws can hold while a further simplification remains impossible. The following model separates \emph{coherent composition} from \emph{a choice of coordinates that removes every phase}.

The coefficient group and the degree of the data matter. The cocycles in Section~\ref{sec:stack} label individual arrows in a view category. Here a phase is attached to a \emph{pair} of group elements. These are different constructions. The finite example is not obtained by renaming the earlier section obstruction as a general second cohomology group.

\subsection{A two-bit phase law}
Write \(\mathbf F_2=\Z/2\Z\) additively, and let
\[
 G=\mathbf F_2^2,\qquad x=(x_1,x_2),\qquad
 c(x,y)=x_2y_1.
\]
All entries and operations in this subsection lie in \(\mathbf F_2\). In particular minus and plus agree, although retaining a minus sign in the coboundary formula helps explain its origin.

\begin{proposition}[A normalized associative phase]\label{KS80}
The phase vanishes when either argument is zero and satisfies
\begin{equation}\label{eq:phase-cocycle}
 c(x,y)+c(x+y,z)=c(y,z)+c(x,y+z).
\end{equation}
There is no function \(a:G\to\mathbf F_2\) such that
\begin{equation}\label{eq:phase-flatten}
 c(x,y)=a(x)+a(y)-a(x+y)\quad\text{for every }x,y\in G.
\end{equation}
\end{proposition}
\begin{proof}
Normalization follows from multiplication by zero. Expanding the two sides of \eqref{eq:phase-cocycle} gives, respectively,
\[
 x_2y_1+(x_2+y_2)z_1,
 \qquad y_2z_1+x_2(y_1+z_1),
\]
which are equal by distributivity.

Every expression \(a(x)+a(y)-a(x+y)\) is symmetric in \(x,y\), because \(G\) is abelian. With \(e_1=(1,0)\) and \(e_2=(0,1)\), however,
\[
 c(e_1,e_2)=0,\qquad c(e_2,e_1)=1.
\]
This contradicts the symmetry required by \eqref{eq:phase-flatten}.
\end{proof}
\source{KS80}
The proof rules out all sixteen possible scalar functions \(a\), not just linear functions or selected trial gauges. The source \code{ProjectiveExample.lean} proves the normalization, cocycle, asymmetry, and no-scalar-flattening statements. The matrix interpretation and group construction below are worked derivations of the finite model, not additional newly compiled declarations.

\plate{30_phase}{The entry in row $x$ and column $y$ is $c(x,y)=x_2y_1$. The two entries singled out in the text, $c(e_1,e_2)=0$ and $c(e_2,e_1)=1$, already obstruct every scalar flattening. Associativity still holds for every triple.}{fig:phase}

\subsection{The phase is visible in actual matrices}
Over \(\R\), define
\[
 X=\begin{pmatrix}0&1\\1&0\end{pmatrix},\qquad
 Z=\begin{pmatrix}1&0\\0&-1\end{pmatrix},\qquad
 U_{(p,q)}=X^pZ^q\quad(p,q\in\{0,1\}).
\]
Direct multiplication gives \(X^2=Z^2=I\) and \(ZX=-XZ\). Moving \(Z^{x_2}\) past \(X^{y_1}\) therefore yields
\begin{equation}\label{eq:phase-matrices}
 U_xU_y=(-1)^{c(x,y)}U_{x+y}.
\end{equation}
For example \(U_{e_1}U_{e_2}=XZ\), while \(U_{e_2}U_{e_1}=-XZ\). The two-bit labels add commutatively, but their chosen matrix lifts do not commute.

Associativity of three lifted matrices requires the same total sign on both parenthesizations. Equation~\eqref{eq:phase-cocycle} is exactly that requirement:
\[
 (U_xU_y)U_z
 =(-1)^{c(x,y)+c(x+y,z)}U_{x+y+z},
\]
\[
 U_x(U_yU_z)
 =(-1)^{c(y,z)+c(x,y+z)}U_{x+y+z}.
\]
Thus a nonzero phase is not automatically a failed associativity law. It can be the data that make the composition law explicit.

Changing lifts to \(U'_x=(-1)^{a(x)}U_x\) changes the phase to
\[
 c'(x,y)=c(x,y)+a(x)+a(y)-a(x+y).
\]
No such change makes \(c'\) vanish identically. The antisymmetric part is unchanged by every scalar change of lift. In this model it is
\[
 \kappa(x,y)=c(x,y)-c(y,x)
            =x_2y_1-y_2x_1
\]
This is a directly observable obstruction certificate for the model. Its nonzero value at \((e_1,e_2)\) is enough, without a search over gauges.

The same symmetry argument rules out removing this sign multiplier by arbitrary nonzero complex scalar lifts as well. A multiplicative scalar coboundary \(b(x)b(y)/b(x+y)\) on an abelian group is symmetric, whereas the multiplier in \eqref{eq:phase-matrices} is not. This last sentence is an algebraic consequence written here. The retained Lean example uses the stated \(\mathbf F_2\)-valued formulation.

\subsection{Retaining the phase restores an ordinary multiplication law}
Instead of trying to discard the phase, include it in the state. On the eight-element set \(\widetilde G=\mathbf F_2\times G\), put
\begin{equation}\label{eq:central-extension}
 (a,x)\cdot(b,y)=(a+b+c(x,y),x+y).
\end{equation}
The first coordinate of the left-associated product of three elements is
\(a+b+d+c(x,y)+c(x+y,z)\). The right-associated coordinate is
\(a+b+d+c(y,z)+c(x,y+z)\). Equation~\eqref{eq:phase-cocycle} proves equality. Normalization gives identity \((0,0)\), and
\[
 (a,x)^{-1}=(a+c(x,x),x).
\]
Both inverse products are the identity because \(x+x=0\) and arithmetic is modulo two. Hence \eqref{eq:central-extension} defines a group.

Projection \((a,x)\mapsto x\) is a homomorphism onto \(G\). Its kernel is \(\{(0,0),(1,0)\}\), which commutes with every element. The map
\[
 (a,x)\longmapsto(-1)^aU_x
\]
is a homomorphism by \eqref{eq:phase-matrices}. Its eight matrices are distinct: they are \(\pm I,\pm X,\pm Z,\pm XZ\). It is therefore faithful. The phase extension is not an assertion that \(G\) itself has acquired an ordinary representation with these same lifts. The represented group has changed.

\begin{insight}
\textbf{A precise form of retaining the obstruction.} The original four labels do not support a phase-free multiplication of the chosen matrices. Keeping one additional bit gives an eight-element multiplication that does. This repairs the description by retaining information, not by proving that the original phase was a coboundary.
\end{insight}

\subsection{How this example fits the realization calculus}
There are three different questions in this small model. First, the phase satisfies the associativity relation. Second, scalar changes cannot trivialize it. Third, an enlarged state that retains the phase admits an ordinary multiplication law. Each question has a complete answer, but their answers are not interchangeable.

The example illustrates why an obstruction should be attached to the allowed changes and to a specified presentation problem. Enlarging the presentation can be mathematically productive, but it changes what is being realized and what information is carried. No general theorem that every obstruction can be removed by a finite enlargement is claimed. The finite checker enumerates all group triples, all scalar gauges, and all multiplication triples in \(\widetilde G\), independently of the symmetry proof above.

\addtocontents{toc}{\protect\newpage}
\input{chapters/05_geometry}
\input{chapters/05a_certificates}
\input{chapters/05b_transport_cost}
\partnote{From approximation to exact realization}{A vanishing defect is a statement about measurements. A limiting realization also requires control of the candidate itself. This part follows the movement of approximate answers, the closure of admissible correction images, and the meaning of the ambient completion.}
\chap{Exactification without a fixed contraction rate}\label{sec:variable-exactification}
The constructed linear repair in Section~\ref{sec:split} is a useful special case of a more general principle. One can prove convergence without one fixed contraction factor, without a linear update rule, and without a globally Lipschitz iteration. What must not be lost is control of the total movement.

\subsection{The two limits that must be established}
Let \((M,\dist)\) be a complete metric space. Let \(\omega:M\to[0,\infty)\) be a nonnegative residual function, and call a point exact when \(\omega(x)=0\). We regard \(\omega\) as a real-valued function when discussing continuity. Assume it is lower semicontinuous, so that every sublevel set \(\{x:\omega(x)\le a\}\) is closed.

For an approximation sequence \(f_n\), the statements
\[
 \omega(f_n)\longrightarrow0
 \quad\text{and}\quad
 f_n\longrightarrow x
\]
are different. The first concerns the observed discrepancy. The second concerns the mathematical candidate in the chosen metric. Neither implies the other without additional hypotheses.

\begin{theorem}[Summable movement gives exactification]\label{KS71}
Suppose \(b_n\ge0\), \(\sum_{n=0}^{\infty}b_n<\infty\), and
\[
 \dist(f_n,f_{n+1})\le b_n,\qquad \omega(f_n)\longrightarrow0.
\]
Then there is a limit \(x\in M\) satisfying
\[
 f_n\longrightarrow x,\qquad\omega(x)=0,\qquad
 \dist(f_n,x)\le\sum_{k=n}^{\infty}b_k.
\]
\end{theorem}
\begin{proof}
For \(m>n\), the triangle inequality gives
\[
 \dist(f_n,f_m)\le\sum_{k=n}^{m-1}b_k.
\]
The tail of a convergent nonnegative series tends to zero, so \((f_n)\) is Cauchy. Completeness supplies its limit \(x\). Taking \(m\to\infty\) in the preceding inequality gives the tail estimate.

For each \(\varepsilon>0\), the residual eventually lies in the closed sublevel set \(\{z:\omega(z)\le\varepsilon\}\). Its limit belongs to that set. Therefore \(\omega(x)\le\varepsilon\) for every positive \(\varepsilon\). Nonnegativity implies \(\omega(x)=0\).
\end{proof}
\source{KS71}
The formal theorem is slightly more general: it is stated in a complete pseudometric space. In a metric space the limit of this fixed sequence is unique. That does not mean different admissible approximation procedures select the same exact point.

\subsection{Variable rates and a movement budget}
Let \(q_n,C_n\ge0\), and define the prefix product
\[
 Q_0=1,\qquad Q_{n+1}=q_nQ_n.
\]
Assume
\begin{equation}\label{eq:variable-rate}
 \omega(f_{n+1})\le q_n\omega(f_n),\qquad
 \dist(f_n,f_{n+1})\le C_n\omega(f_n).
\end{equation}

\begin{theorem}[Variable-rate exactification]\label{KS72}
If \(Q_n\to0\) and \(\sum_n C_nQ_n<\infty\), then \((f_n)\) has an exact limit and
\begin{equation}\label{eq:variable-tail}
 \dist(f_n,x)\le\omega(f_0)\sum_{k=n}^{\infty}C_kQ_k.
\end{equation}
For constant \(q_n=q\in[0,1)\) and constant \(C_n=C\), the sharper restarted estimate is
\[
 \dist(f_n,x)\le\frac{C}{1-q}\,\omega(f_n).
\]
\end{theorem}
\begin{proof}
Induction in \eqref{eq:variable-rate} yields \(0\le\omega(f_n)\le Q_n\omega(f_0)\). The residual therefore tends to zero, and the step sizes are bounded by the summable sequence \(C_nQ_n\omega(f_0)\). Apply Theorem~\ref{KS71}.

For constant rates, restart the same induction at index \(n\). The residual at \(n+k\) is at most \(q^k\omega(f_n)\), so the movement from \(f_n\) to its limit is at most \(C\omega(f_n)\sum_{k\ge0}q^k\). Summing the geometric series proves the last estimate. No division by an individual residual is needed, so a zero residual causes no exceptional case.
\end{proof}
\source{KS72}
The individual factors \(q_n\) need not all be below one. Temporary residual increases are permitted if the prefix products decay and the weighted movements remain summable. Conversely, a decaying product is not enough when its movement budget diverges.

\subsection{A convergent path with no uniform geometric factor}
In \(\R^2\) with the maximum norm, fix \(s\in\R\) and take
\[
 f_n=(s,(n+1)^{-1}),\qquad \omega(u,r)=|r|.
\]
The protected coordinate is the first one. Direct calculation gives
\[
 q_n=\frac{n+1}{n+2},\quad Q_n=\frac1{n+1},\quad
 C_n=\frac1{n+2},\quad
 \dist(f_n,f_{n+1})=\frac1{(n+1)(n+2)}.
\]
No constant \(q<1\) bounds all the ratios \(q_n\), since they tend to one. Nevertheless
\[
 \sum_{k=n}^{\infty}C_kQ_k
 =\sum_{k=n}^{\infty}\left(\frac1{k+1}-\frac1{k+2}\right)
 =\frac1{n+1}.
\]
The limit is \((s,0)\), and the error bound is an equality. The model shows why the variable-rate theorem is not merely a rephrasing of a fixed geometric contraction.

\plate{25_summable_path}{Polynomial residual decay can still give an exact limit with a sharp error certificate. The residual $1/(n+1)$ is not summable, but the actual movement $1/((n+1)(n+2))$ is. The theorem uses the movement series.}{fig:summable-path}

\subsection{Residual bursts can be compatible with convergence}
Set \(q_{2k}=2\), \(q_{2k+1}=1/8\). Then
\[
 Q_{2k}=4^{-k},\qquad Q_{2k+1}=2\cdot4^{-k}.
\]
For \(f_n=(s,Q_n)\), the residual doubles at the start of every pair of steps. Choose \(C_{2k}=1\) and \(C_{2k+1}=7/8\). These give equality in the movement bound, and
\[
 \sum_{n=0}^{\infty}C_nQ_n
 =\sum_{k=0}^{\infty}\left(1+\frac74\right)4^{-k}
 =\frac{11}{3}.
\]
Thus \(f_n\to(s,0)\) despite infinitely many increases. Starting at an even index \(2k\), the remaining movement budget is \(\frac{11}{3}4^{-k}\). At an odd index \(2k+1\), it is \(\frac83 4^{-k}\). These bounds are larger than the actual distance because the path repeatedly moves away and returns.

\plate{26_residual_bursts}{The residual increases at every even-to-odd step, but each two-step product is $1/4$. A convergent product and a summable movement budget control the whole path. Monotonicity of individual residual values is not assumed by the variable-rate theorem.}{fig:residual-bursts}

\subsection{Protected meaning survives a controlled limit}
Return to Banach and normed linear spaces, with continuous linear maps \(d:E\to R\) and \(p:E\to S\). Put \(\omega(f)=\norm{df}\). The residual is continuous, hence lower semicontinuous.

\begin{proposition}[Controlled paths kill the actual obstruction]\label{KS73}
Suppose \(E\) is complete, \(\norm{df_n}\to0\), \(p(f_n)=p(f_0)\), and \(\dist(f_n,f_{n+1})\le b_n\) for a summable nonnegative budget. Then there is \(x\) with
\[
 f_n\to x,\qquad dx=0,\qquad p(x)=p(f_0),
\]
and the constrained obstruction \([df_0]\in R/d(\ker p)\) vanishes.
\end{proposition}
\begin{proof}
Theorem~\ref{KS71} gives the limit and \(dx=0\). By continuity, \(p(x)=\lim_n p(f_n)=p(f_0)\). Consequently \(a=f_0-x\) lies in \(\ker p\) and satisfies \(da=df_0\). It is an actual admissible correction, so the quotient class is zero.
\end{proof}
\source{KS73}
This proposition converts a controlled approximation argument into a statement about the algebraic quotient defined earlier. Merely showing that \(df_n\to0\) would not produce the admissible vector \(f_0-x\), because there might be no candidate limit \(x\).

\subsection{When two repair paths select the same point}
\begin{proposition}[Path independence requires a comparison]\label{KS74}
Let \(f_n\to x\) and \(g_n\to y\) in a metric space. If
\[
 \dist(f_n,g_n)\le\beta_n,\qquad\beta_n\to0,
\]
then \(x=y\). In contrast, exactness and preservation of a nonzero observable alone do not force a unique answer.
\end{proposition}
\begin{proof}
Continuity of distance and the squeeze theorem give \(\dist(x,y)=0\), hence equality. For the last assertion use \(E=\R\times(\R\times\R)\), defect \(d(s,u,r)=r\), and observable \(p(s,u,r)=s\). The points \((1,0,0)\) and \((1,1,0)\) are distinct exact answers with the same nonzero protected value.
\end{proof}
\source{KS74}
A stronger worked example uses actual convergent paths from the \emph{same} initial input:
\[
 f_n^{(a)}=(s,a(1-2^{-n}),2^{-n})\qquad(a\in\R).
\]
All paths begin at \((s,0,1)\), keep \(p=s\), and have residual \(2^{-n}\). Their maximum-norm movements are \(\frac12\max\{|a|,1\}2^{-n}\), so every path satisfies geometric exactification. Their limits are \((s,a,0)\). For different \(a,b\),
\[
 \dist(f_n^{(a)},f_n^{(b)})=|a-b|(1-2^{-n})\longrightarrow|a-b|,
\]
not zero. The missing uniqueness hypothesis is visible in the exact direction along which the paths separate.

\plate{27_path_limits}{All paths start at the same candidate, have the same residual at each step, and preserve the same coordinate $s$. They converge to different exact points because the coordinate $u$ is unprotected. A vanishing residual is not a path-comparison budget.}{fig:path-limits}

\subsection{Changing the sampling schedule does not preserve a rate}
A residual profile \(\eta_n\) indexed by resolution is not an iteration trajectory. Nevertheless a useful convergence property is invariant under a cofinal schedule when the profile is monotone.

\begin{proposition}[Cofinal sampling of a decreasing profile]\label{KS75}
Suppose \(\eta_n\ge0\) is nonincreasing and \(\phi:\N\to\N\) is monotone and unbounded. Then
\[
 \eta_n\to0\quad\Longleftrightarrow\quad\eta_{\phi(n)}\to0.
\]
Without the monotonicity of \(\eta\), the reverse implication can fail even for a cofinal schedule.
\end{proposition}
\begin{proof}
The schedule tends to infinity, so convergence of the full profile implies convergence of the sampled one. Conversely, choose a sampled value below any prescribed \(\varepsilon>0\). Every later index of the original profile has no larger value, by monotonicity. This proves convergence of the full profile.

For failure without monotonicity, take \(\eta_{2n}=0\), \(\eta_{2n+1}=1\), and \(\phi(n)=2n\). All sampled values vanish, but infinitely many original values equal one.
\end{proof}
\source{KS75}
Even in the decreasing case, the numerical rate changes. The profile \(\eta_n=1/(n+1)\), sampled at \(\phi(n)=2^n-1\), becomes \(2^{-n}\). This is exponential decay in the \emph{number of observations}, not a proof of a faster correction process at the original resolution scale.

\subsection{Three hypotheses that cannot simply be dropped}
\textbf{Summable movement.} In maximum-norm \(\R^3\), let
\[
 f_n=(s,H_n,(n+1)^{-1}),\qquad
 H_n=\sum_{j=1}^{n}\frac1j.
\]
The first coordinate is protected and the residual is the last coordinate. It tends to zero with the same ratios as the polynomial example. The movement equals \(1/(n+1)\), however, and the unprotected coordinate diverges. Here \(C_n=1\), so \(\sum C_nQ_n\) is the divergent harmonic series. Residual decay and semantic preservation hold, but exactification of this path fails.

\textbf{Completeness.} In the metric space \((0,\infty)\), take \(f_n=1/(n+1)\) and \(\omega(x)=x\). The steps are summable and the residual tends to zero, but there is no limit in this space. Completing the space adds the missing zero. It is essential to say whether that completion belongs to the intended model.

\textbf{Lower semicontinuity.} In \(\R\), define \(\omega(0)=1\) and \(\omega(x)=0\) for \(x\ne0\). The same positive sequence converges with summable movement and residual zero at every step, yet its limit is not exact. The zero sublevel set is not closed. This example isolates why convergence of measurements along a sequence must interact correctly with the ambient topology.

\clearpage
\chap{The residual floor and the closure of admissible changes}\label{sec:floor}
An algebraic obstruction asks whether a residual is exactly the defect of an allowed correction. A numerical floor asks how closely allowed defects can approximate it. These questions agree under a closedness hypothesis. Without that hypothesis they can disagree even in a complete Hilbert setting.

\subsection{Distance to the actual correction image}
Let \(R\) be a normed additive group and \(B\le R\) an additive subgroup. In the protected linear model, \(B=d(\ker p)\). For a residual \(\Omega\in R\), define
\[
 \nu_B(\Omega)=\inf_{b\in B}\norm{\Omega-b}.
\]
The subgroup contains zero, so this is a finite nonnegative real number. It is an infimum, not an assertion that a nearest admissible residual exists.

\begin{proposition}[The floor detects closure]\label{KS76}
For every \(b\in B\),
\[
 0\le\nu_B(\Omega)\le\norm{\Omega-b}.
\]
Moreover,
\[
 \nu_B(\Omega)=0\quad\Longleftrightarrow\quad\Omega\in\overline B.
\]
If \(B\) is closed, this is equivalent to \(\Omega\in B\).
\end{proposition}
\begin{proof}
Nonnegativity and the upper bound follow from the definition of an infimum. If the infimum is zero, every positive-radius ball around \(\Omega\) meets \(B\), so \(\Omega\in\overline B\). Conversely, membership in the closure supplies, for every \(\varepsilon>0\), a point \(b\in B\) with \(\norm{\Omega-b}<\varepsilon\). The nonnegative infimum must be zero. Closedness replaces \(\overline B\) by \(B\).
\end{proof}
\source{KS76}
Translation by a member \(b_0\in B\) does not change the floor: replacing \(b\) by \(b-b_0\) gives the same set of distances. Thus the floor is well defined on the algebraic quotient \(R/B\). When \(B\) is a linear subspace, this is the usual quotient seminorm. A nonzero algebraic class can have seminorm zero when \(B\) is not closed. The retained source implements the floor on residual representatives and proves the closure criterion. This quotient interpretation is its elementary translation-invariance calculation, not a new claim that every obstruction quotient carries a separated norm.

\subsection{A precise three-way distinction in this fixed model}
For this specified subgroup and residual there are exactly three possibilities:
\begin{center}
\begin{tabularx}{\linewidth}{@{}p{.29\linewidth}p{.20\linewidth}X@{}}
\toprule Residual location & Algebraic class & Meaning of the floor\\\midrule
\(\Omega\in B\) & Zero & Exact cancellation exists, and the floor is zero.\\
\(\Omega\in\overline B\setminus B\) & Nonzero & Arbitrarily small corrected residuals exist, but no exact cancellation exists.\\
\(\Omega\notin\overline B\) & Nonzero & The floor is strictly positive and blocks every allowed correction.\\\bottomrule
\end{tabularx}
\end{center}
This is a classification for the stated normed-group correction problem, not a universal trichotomy for every possible Komorebi model. Testing which case holds can itself be difficult. The definitions do not supply a decision algorithm.

In finite-dimensional real or complex linear models, every linear subspace is closed. The middle row is therefore absent. It becomes possible in infinite dimensions, or for additive subgroups that are not closed real-linear subspaces. For example the rational subgroup of \(\R\) is dense, so the distance from any irrational real to that subgroup is zero even though its coset is nonzero.

Closedness also should not be confused with a numerical error bound back in the candidate space. The statement \(\Omega\in B\) supplies an admissible correction, but not a bound on its size. The closed-range estimates of Sections~\ref{sec:closed} and~\ref{sec:motion}, or their protected counterparts, are the additional analytic tools that control that size in their declared settings.

\subsection{A persistent floor blocks every correction sequence}
The correction image can vary with the resolution or the admissibility rule. Let \(B_n\le R\) and let \(\Omega_n\in R\) be given.

\begin{proposition}[Uniformly positive floors obstruct asymptotic repair]\label{KS77}
Suppose \(c>0\) and eventually
\[
 c\le\nu_{B_n}(\Omega_n).
\]
For any choices \(b_n\in B_n\), it is impossible that
\[
 \norm{\Omega_n-b_n}\longrightarrow0.
\]
\end{proposition}
\begin{proof}
Every allowed correction satisfies
\(\norm{\Omega_n-b_n}\ge\nu_{B_n}(\Omega_n)\). The corrected norms are therefore eventually at least \(c\). They cannot converge to zero.
\end{proof}
\source{KS77}
The quantifier covers all admissible sequences, not just the output of one repair routine. On the other hand, the hypothesis is an eventually \emph{uniform} positive lower bound. Having a nonzero obstruction class at each stage is not the same statement. The classes might have zero floors, or their positive floors might tend to zero.

A finite-dimensional illustration uses \(R=\R^2\), \(B=\R\times\{0\}\), and \(\Omega_n=(n,1)\). Every correction leaves the second coordinate equal to one, so the floor is one independently of \(n\). Replacing the residuals by \((n,1/(n+1))\) gives nonzero algebraic classes at every stage, but their floors tend to zero. The persistent-floor theorem applies to the first sequence and not the second.

\subsection{The bridge back to exactification}
Theorem~\ref{KS73} and Proposition~\ref{KS76} fit together closely. If corrections produce a norm-convergent, protected path with vanishing residual, the actual quotient class vanishes. If only the corrected residuals tend to zero, all that follows is membership in \(\overline{d(\ker p)}\).

In particular, when the class is nonzero but its floor is zero, \emph{no} protected path with vanishing residual can satisfy all the exactification hypotheses. Something must fail: the path may escape, its movements may not be summable, the candidate space may not be complete, or its residual and semantic maps may not have the required continuity. The next model isolates escape while keeping the spaces Hilbert and both maps continuous.

\clearpage
\chap{Exact finite repairs can escape the intended ambient space}\label{sec:escape}
This model links the algebraic obstruction, the numerical floor, the movement of approximations, and the semantic image of an inverse limit. It uses a diagonal operator whose entries are completely explicit. Every finite problem is solvable. The protected infinite problem is not.

\subsection{The complete spaces and the preservation rule}
Index sequences by the positive integers. Set
\[
 E=\ell^2(\N_{>0};\R)\oplus\R,\qquad
 R=\ell^2(\N_{>0};\R),\qquad
 \norm{(u,t)}_E^2=\sum_{n\ge1}|u_n|^2+|t|^2.
\]
Let \(w=(1/n)_{n\ge1}\). It belongs to \(\ell^2\), because
\[
 \sum_{n\ge1}\frac1{n^2}
 \le1+\sum_{n\ge2}\frac1{n(n-1)}=2.
\]
Define the defect and protected observable by
\[
 d(u,t)_n=\frac{u_n+t}{n},\qquad p(u,t)=t.
\]
The diagonal map \(T:u\mapsto(u_n/n)_n\) has operator norm one. Since \(d(u,t)=Tu+tw\),
\[
 \norm{d(u,t)}\le\norm u+|t|\norm w
 \le\sqrt{1+\norm w^2}\,\norm{(u,t)}_E.
\]
Thus both \(d\) and \(p\) are continuous linear maps between the declared complete spaces. No failure of completeness or continuity is built into the model.

Take the raw candidate \(f_0=(0,1)\). Its residual is \(w\), and the coordinate that must be preserved is \(t=1\). An exact protected answer would require
\[
 u_n=-1\quad\text{for every }n.
\]
That sequence does not lie in \(\ell^2\). Consequently no answer in \(E\) is both exact and protected.

In fact \(\ker d=\{0\}\). If \(d(u,t)=0\), then \(u_n=-t\) for every \(n\). A constant square-summable sequence must have constant zero, so \(t=0\) and \(u=0\). There is an exact answer after abandoning the preservation rule, namely zero, but it changes \(p(f_0)=1\) to zero. Even a trivial ambiguity space does not imply feasibility.

\subsection{The obstruction class is nonzero while its floor is zero}
The admissible correction space is \(\ker p=\ell^2\oplus\{0\}\), with correction image
\[
 B=d(\ker p)=\ran T
  =\{r\in\ell^2:(nr_n)_{n\ge1}\in\ell^2\}.
\]
It contains every finitely supported sequence, so it is dense in \(R\). But \(w\notin B\), because \((nw_n)_n=(1,1,\ldots)\) is not square summable. Therefore
\[
 [w]\ne0\text{ in }R/B,\qquad \nu_B(w)=0.
\]
This is the middle row of the preceding classification. The example is a worked Hilbert-space derivation and is not a newly compiled Lean instance. Its finite-coordinate calculations are independently tested. The general floor and exactification results are linked by KS71--KS77 to the retained formal source.

\subsection{A complete calculation of the finite repairs}
Let \(\mathbf 1_N\) be one in coordinates \(1,\ldots,N\) and zero elsewhere. Put
\[
 a_N=(\mathbf 1_N,0),\qquad f_N=f_0-a_N=(-\mathbf 1_N,1).
\]
Then \(p(f_N)=1\), and the first \(N\) residual coordinates vanish exactly. The remaining residual and movement satisfy
\[
 d(f_N)_n=\begin{cases}0&n\le N,\\1/n&n>N,\end{cases}
 \qquad \norm{a_N}^2=N,\qquad \norm{f_N}^2=N+1.
\]
For every \(N\ge1\), telescoping comparisons give
\begin{equation}\label{eq:escape-tail}
 \frac1{N+1}\le\norm{d(f_N)}^2
 =\sum_{n>N}\frac1{n^2}\le\frac1N.
\end{equation}
Indeed \(1/[n(n+1)]\le1/n^2\le1/[n(n-1)]\) for \(n\ge2\). Summing the two telescoping bounds proves \eqref{eq:escape-tail} without numerically evaluating an infinite series.

The residual tends to zero, but the corrections grow like \(\sqrt N\). Successive candidates differ by a unit coordinate, so
\[
 \norm{f_{N+1}-f_N}=1.
\]
The movement series diverges and the sequence is not Cauchy. This is exactly the hypothesis missing from the summable exactification theorem.

\begin{center}
\begin{tabular}{@{}rrrr@{}}
\toprule Repaired coordinates \(N\) & \(\norm{a_N}\) & Residual square interval & Finite inverse gain\\\midrule
1 & 1 & \([1/2,1]\) & 1\\
4 & 2 & \([1/5,1/4]\) & 4\\
16 & 4 & \([1/17,1/16]\) & 16\\
64 & 8 & \([1/65,1/64]\) & 64\\\bottomrule
\end{tabular}
\end{center}
The intervals bound the \emph{infinite} residual remaining after the displayed repair. The inverse-gain column concerns a different finite operator, defined below. Keeping those meanings separate prevents a finite success from being mistaken for a uniform infinite certificate.

\plate{28_escaping_repair}{The squared correction size is $N$, while the infinite residual square lies between $1/(N+1)$ and $1/N$. The plotted band consists of proved bounds, not a numerical truncation of the tail. Smaller residuals require larger changes and do not produce a convergent candidate.}{fig:escaping-repair}

\subsection{Escape is not an accident of this particular algorithm}
Suppose an arbitrary protected candidate \((u,1)\) satisfies \(\norm{d(u,1)}\le\varepsilon\). On its first \(N\) coordinates,
\[
 \sum_{n=1}^{N}|u_n+1|^2
 \le N^2\sum_{n=1}^{N}\frac{|u_n+1|^2}{n^2}
 \le N^2\varepsilon^2.
\]
The reverse triangle inequality in \(\R^N\) gives
\begin{equation}\label{eq:escape-necessity}
 \norm u\ge\sqrt N-N\varepsilon\qquad(N\ge1).
\end{equation}
If a sequence of protected candidates has residual tending to zero, then for each fixed \(N\), its eventual correction norms are at least, for example, \(\sqrt N/2\). Taking \(N\) arbitrarily large shows that these norms must tend to infinity. There is no uniformly bounded family of protected corrections with arbitrarily small residual.

Thus changing the iteration rule cannot fix the obstruction while preserving all the declared requirements. One must change the input, the preservation contract, or the intended ambient class. The inequality is a quantitative obstruction to a bounded repair procedure, rather than a failure report from one chosen recurrence.

\subsection{Every finite operator is invertible, but the gains diverge}
On \(\R^N\), let
\[
 T_N=\operatorname{diag}(1,1/2,\ldots,1/N).
\]
It is invertible and \(\norm{T_N^{-1}}=N\) in the Euclidean norm. The upper bound follows by multiplying each coordinate by an integer at most \(N\), and equality is attained at the last coordinate vector. Full finite rank therefore supplies no uniform gain as \(N\to\infty\).

In the infinite space, \(\norm{e_n}=1\) but \(\norm{Te_n}=1/n\), while \(\ker T=\{0\}\). A uniform estimate
\(\dist(u,\ker T)\le C\norm{Tu}\) would force \(1\le C/n\) for every \(n\), which is impossible. This agrees with the nonclosed range already established directly. The low residual directions do not reflect new exact vectors. They reflect increasingly ill-conditioned approximate directions.

\subsection{A coherent inverse-limit point still exists in a larger space}
At each finite level, the protected exact equations have the unique solution
\[
 (-1,\ldots,-1;1)\in\R^N\oplus\R.
\]
Truncation maps the solution at level \(N+1\) to the one at level \(N\). These finite solution sets are singletons, so their inverse limit is also a singleton. Its coordinate description is
\[
 (-1,-1,\ldots;1)\in\R^{\N_{>0}}\times\R.
\]
The coherent point is not missing. It lies outside the intended \(\ell^2\)-based semantic space.

This example does not contradict finite-fiber existence. It distinguishes the ordinary inverse limit of finite views from the image of a separate comparison map that sends square-summable semantic objects to their finite prefixes. The prescribed coherent point is not in that image. Replacing \(\ell^2\) by all sequences makes the equation solvable, but changes the topology and the class of allowed objects.

\begin{insight}
\textbf{Four statements, four answers.} Every finite prefix is exactly solvable. The ordinary coherent inverse-limit point exists. Its protected Hilbert realization does not exist. Approximate protected Hilbert realizations have residuals tending to zero only at the cost of unbounded correction size. All four facts hold simultaneously in one explicit model.
\end{insight}

\input{chapters/06_cases}
\chap{Mathematical meaning and the remaining boundaries}\label{sec:meaning}
\subsection{A relative theory can be complete without being absolute}
The strict space \(\Gamma(F)\) is canonical relative to its diagram: its universal property determines it up to the appropriate unique isomorphism. The additive obstruction is canonical relative to its admissible corrections. A marked realization problem is invariant under the specified comparison equivalences. A repair can be canonical relative to a protected observable, or orthogonal relative to a fixed inner product.

These are meaningful and different kinds of canonicity. They need not culminate in one object independent of every choice of views, coefficients, topology, semantics, and geometry. Some changes of view intentionally forget information. Others strengthen the question. Treating every such change as an equivalence would remove the very distinctions the calculus is designed to preserve.

The appropriate central statement is therefore not that every possible mathematical presentation lives in one fixed universal space. It is that a declared system of representations can be organized into coherent realizations, with explicit extension, comparison, and repair problems. A proposed bridge between two such systems has to show what structure it preserves and what information it loses.

\subsection{Internal realization and intended objects}
Suppose a separate category \(\mathcal S\) contains the intended semantic objects of an application, and a realization construction gives a comparison functor
\[
 \mathcal R:\mathcal S\longrightarrow\mathcal K,
\]
where \(\mathcal K\) is the relevant category or groupoid of coherent data. Defining \(\mathcal K\) does not establish that \(\mathcal R\) is faithful, full, or essentially surjective. These are new mathematical questions about the actual semantic comparison.

Faithfulness asks whether distinct semantic arrows can become indistinguishable. Fullness asks whether every compatible comparison comes from a semantic one. Essential surjectivity asks whether every intended coherent datum is realized by an actual semantic object. Depending on the application, one may seek only a restricted image or a classified obstruction to belonging to it.

This distinction prevents a formal globalization of encoded data from being mistaken for an unrestricted arithmetic or physical existence theorem. The abstract construction does not manufacture missing semantic objects.

\subsection{Every finite test can pass while the infinite problem fails}
\begin{example}[Finite prefixes without a global point]\label{KS54}
Let \(X_n=\{m\in\N:m\ge n\}\), with bonding maps given by inclusion. Every level is nonempty, and every finite prefix has a coherent point. The inverse limit is nevertheless empty.
\end{example}
\begin{proof}
For a prefix ending at \(N\), assign the integer \(N\) at every level. All inclusions preserve it. In a global point, however, the inclusion equations force every coordinate to be the same integer \(m\), while membership at level \(n\) requires \(m\ge n\) for every \(n\). Taking \(n=m+1\) is a contradiction.
\end{proof}
\source{KS54}
This counterexample is formalized in \code{FinitePrefixBoundary.lean} using the original resolution and partial-section types. It explains why inverse-limit existence needs more than separate finite solvability. Finiteness, compactness, surjectivity, or another coherence principle must be proved in the relevant application. The example does not contradict the strict gluing theorem: a compatible matching family of all prefixes was never supplied.

\subsection{A positive infinite-existence criterion}
The failed example isolates a missing compactness or lifting hypothesis. Finite fibers supply a concrete positive criterion without requiring every bonding map to be surjective.

\begin{proposition}[Finite fibers and persistent images]\label{KS58}
Let \((X_n,\rho_n)\) be a sequential inverse system of nonempty finite sets. Its inverse limit is nonempty. More precisely, if \(\rho_{n,m}:X_m\to X_n\) is composite restriction, then
\[
 \im(\operatorname{ev}_n)=Y_n,
 \qquad Y_n=\bigcap_{m\ge n}\im\rho_{n,m}.
\]
Thus a prescribed value at level \(n\) extends precisely when it survives every later image.
\end{proposition}
\begin{proof}
The images in the finite set \(X_n\) form a decreasing sequence of nonempty subsets, so they stabilize to a nonempty \(Y_n\). The maps \(\rho_n\) send \(Y_{n+1}\) into \(Y_n\). They are onto these subsets: given \(y\in Y_n\), choose a later stage \(m\) after both image sequences at levels \(n\) and \(n+1\) have stabilized. Lift \(y\) to \(X_m\) and restrict that lift to \(X_{n+1}\). It belongs to \(Y_{n+1}\) and maps to \(y\).

Choose successive preimages in these nonempty finite fibers. They give a coherent family. Starting at any specified \(y\in Y_n\), and restricting to earlier levels, proves the claimed evaluation image. The reverse inclusion follows by restricting a global family from each later stage.
\end{proof}
\source{KS58}
This is the standard finite-set inverse-limit principle, here with its sequential persistent-image proof~\cite{stacks-limits}. Its formalization in \code{FiniteRealizationExistence.lean} constructs an explicit equivalence with the pinned library's section type, then applies the finite cofiltered-system and eventual-range theorems. The printed sequential proof and this formal route establish the same stated criterion. Both use ordinary classical choice. Surjective bonding maps on arbitrary nonempty sets give a second sufficient condition by direct recursive lifting. Neither criterion says that arbitrary data in a non-surjective finite tower extend.

\subsection{A finite computation that exposes the preservation contract}
Consider a raw vector \(x=(2,0,5)\in\R^3\) with defect \(d(x)=x_1-x_2\). Unconstrained orthogonal repair gives \((1,1,5)\), with squared error two. If the first and third coordinates must be preserved, the unique exact repair is \((2,2,5)\), with squared error four.

Both outputs are exact. Only one satisfies the second preservation problem. Neither theorem contradicts the other because the admissible changes differ. An implementation that reports only zero residual would not distinguish the two tasks.

For a finite linear system over a field, the same question admits independently checkable positive and negative certificates. Let \(A\in\K^{m\times n}\) encode compatibility, \(P\in\K^{p\times n}\) the protected values, and \(x\in\K^n\) the raw candidate. Let the columns of \(N\) be a basis of \(\ker P\). The correction equation is
\[
 ANz=Ax,\qquad y=x-Nz.
\]
It is equivalent to the stacked system \(Ay=0\), \(Py=Px\).

\begin{proposition}[A finite certificate alternative]\label{KS57}
Exactly one of the following occurs:
\begin{enumerate}
\item There is \(y\) with \(Ay=0\) and \(Py=Px\). Every answer is \(y+H\), where \(H=\ker A\cap\ker P\).
\item There are row vectors \(\lambda\in\K^{1\times m}\) and \(\mu\in\K^{1\times p}\) with
\[
 \lambda A=\mu P,\qquad\lambda Ax=1.
\]
\end{enumerate}
The negative witness annihilates every admissible correction defect, since \(\lambda AN=\mu PN=0\). The ambiguity dimension in the positive case is
\[
 \dim H=n-\operatorname{rank}\begin{pmatrix}A\\P\end{pmatrix}.
\]
\end{proposition}
\begin{proof}
Row-reduce \(T=\begin{pmatrix}A\\P\end{pmatrix}\) with right side \(b=\begin{pmatrix}0\\Px\end{pmatrix}\), retaining the invertible row-operation matrix. A consistent reduced system supplies an answer, and differences of answers are exactly \(\ker T\). Rank-nullity gives the dimension.

Otherwise a row combination \(w=(u,v)\) satisfies \(wT=0\) and \(wb=\alpha\ne0\). Put \(\lambda=-u/\alpha\), \(\mu=v/\alpha\). Then \(\lambda A=\mu P\) and \(\lambda Ax=\mu Px=1\). Conversely an exact protected answer would imply
\[
 1=\lambda Ax=\mu Px=\mu Py=\lambda Ay=0,
\]
which is impossible in a field. Therefore the alternatives are exclusive and exhaustive.
\end{proof}
\source{KS57}
The formal proof gives these literal row-vector equations over an arbitrary field. It uses algebraic dual separation of the stacked image rather than verifying the row-reduction program. The affine answer equivalence and rank-nullity formula are proved separately. Thus a normalized obstruction is an actual checkable witness, while universal correctness of the supplied Python elimination implementation remains a different question.

\subsection{One equation, three preservation problems}
For the following three cases work over \(\Q\), with \(A=(1,-1,0)\) and \(x=(2,0,5)\). With \(P(x)=(x_1,x_3)\), the unique answer is \((2,2,5)\). With \(P(x)=x_3\), the entire answer set is \(\{(t,t,5):t\in\Q\}\). With \(P(x)=(x_1,x_2)\), no answer exists. In this last case
\[
 \lambda=(1/2),\qquad\mu=(1/2,-1/2)
\]
is a normalized negative certificate over \(\Q\) or \(\R\): \(\lambda A=\mu P\) and \(\lambda Ax=1\).

The package implements this finite interface over \(\Q\) in \code{examples/protected\_repair.py}. It returns a particular exact answer and an ambiguity basis, or a normalized dual witness. Its verifier checks the certificate equations directly. A Gaussian particular answer is \emph{not} advertised as the nearest answer of Theorem~\ref{KS56}. That is a separate geometric selection.

Exact field elimination does not solve the integer parity problem by changing its coefficient domain. Over \(\Q\), \(2m=1\) has the solution \(1/2\). Over \(\Z\), it has none. The rational interface rejects floating-point inputs rather than silently treating an approximate decimal as exact arithmetic. Its finite tests establish the tested implementation behavior, not an arbitrary-input Lean refinement theorem.

\subsection{Why this is relevant to machine reasoning}
The framework gives a machine explicit mathematical objects to manipulate: views, restrictions, partial realizations, typed defects, admissible corrections, and comparison witnesses. A successful step can carry both an output and a certificate of what remained fixed. A failed step can return an obstruction in a specified quotient rather than an unexplained failure flag.

The split coordinate theorem is especially direct. In its stated class, the same candidate is represented by protected information and residual data, and repair zeros only the residual coordinate. The descent theorem explains which local certificates can be assembled. The analytic results explain when a small measured defect has a quantitatively reliable meaning.

These observations motivate a machine oriented mathematical language. They do not establish an autonomous reasoning system, an empirical intelligence gain, or a general efficient algorithm. Classical choice, infinite products, quotient membership, and bounded inverse constructions can be mathematically valid without yielding executable finite procedures. Implementation claims need their own specifications and tests.

\subsection{What this monograph closes, and what it leaves open}
The book supplies a coherent account of strict realization spaces, additive and nonabelian extension problems, split tower repair, retained multiscale detail, a genuine composition-site stack, presentation transport, graph holonomy, closed-range stability, and moving exact spaces. The additional convergence and obstruction development separates small measurements from small movements, coherent coordinate families from their semantic image, and zero residual from zero transformation cost. Each layer is attached to the data that make it meaningful.

The following broader targets are not conclusions of this edition: arbitrary higher categorical coefficients and infinity-stack descent, unrestricted changes of view doctrine, norm reconstruction for every bounded infinite residual hierarchy, nonlinear or unknown-connection repair, perturbation across changing Hilbert geometries without new comparison data, and semantic essential surjectivity for unrestricted arithmetic or physical applications.

The integrated proof infrastructure for the stated source closure is supplied in this edition. The canonical project has \EditionFormalModules\ production modules, a single root, clean compiler records, declaration axiom reports, and rejection controls. All five previously identified written-only items have formal counterparts. Further historical material is not silently included in this closure. Extensions of the theory require new hypotheses, sources, and verification records.

\begin{insight}
\textbf{The organizing principle.} Coherent realization, obstruction to extension, freedom of completion, and stability of repair belong to one mathematical conversation. Their unification is useful precisely when it does not identify them with one another.
\end{insight}

\appendix
\clearpage
\chap{A reader's workshop}\label{app:workshop}
These short problems isolate choices that recur in the main development. Each has a complete solution. They are reading exercises and direct applications of the cited theory, not additional entries in the formal theorem inventory.

\subsection*{1. A stable limit does not make every seed persistent}
At every level take \(X_n=\{0,1\}\), with each bonding map sending both elements to zero. Find the inverse limit and the image of evaluation at every level. Does nonemptiness of every finite fiber make the seed one extendable?

\textbf{Solution.} For a coherent family, \(x_n=\rho_n(x_{n+1})=0\) for every \(n\). The all-zero family is coherent and is therefore the only point. Every evaluation image is \(\{0\}\). Equivalently, the image of restriction from any strictly later level is already \(\{0\}\), so the persistent image is \(\{0\}\). The seed one never extends. This is compatible with KS58, which proves nonemptiness and characterizes the persistent image, not surjectivity of every evaluation.

\subsection*{2. A residual chart can change without changing the defect}
Let \(d:\R^2\to\R\) be \(d(u,r)=r\). For \(a\in\R\), let \(A_a(r)=(ar,r)\). Compute the exact repair \(P_a\), compare it with \(P_b\), and decide when they commute.

\textbf{Solution.} Since \(dA_a=I\),
\[
 P_a(u,r)=(u-ar,0),\qquad
 P_b(u,r)-P_a(u,r)=((a-b)r,0).
\]
Both fix the horizontal exact axis. Thus \(P_aP_b=P_b\) and \(P_bP_a=P_a\). They commute exactly when they are equal, which occurs exactly when \(a=b\), as can be checked at input \((0,1)\). Every repair has zero residual. The different answers reflect the different right inverses, as in KS65.

\subsection*{3. Pairwise invertibility does not supply a triple law}
Let \(t\) be a nonidentity transposition in \(S_3\). Suppose three proposed overlap comparisons require frames \(b_0,b_1,b_2\) satisfying
\[
 b_1=tb_0,\qquad b_2=tb_1,\qquad b_2=tb_0.
\]
Every comparison is individually invertible. Can the three frames exist?

\textbf{Solution.} The first two equations give \(b_2=t^2b_0=b_0\). The third then gives \(b_0=tb_0\), and right cancellation gives \(t=1\), contrary to the assumption. Invertibility of every proposed overlap map does not prove the composition law needed for descent. KS34 applies to data satisfying that law, not to arbitrary lists of pairwise isomorphisms.

\subsection*{4. Zero noisy output is not a certificate of reference exactness}
Take \(D(a,b)=a\), \(\widetilde D(a,b)=a+\tfrac1{10}b\), and input \(x=(-10,100)\). The measured operator returns zero exactly. Find the distance to the reference kernel and evaluate the reference uncertainty bound.

\textbf{Solution.} The reference kernel is the vertical axis, so the distance is ten. Its residual gain is one, and the operator error norm is \(1/10\). With \(y=0\) and \(\eta=0\), KS69 gives
\[
 \dist(x,\ker D)\le\frac1{10}\sqrt{10^2+100^2}=\sqrt{101}.
\]
The true distance ten is below this bound. A formula omitting the operator-error term would instead claim a zero distance, which is false. The input is exact for the measured model and not for the reference one.

\subsection*{5. A prefix computation can hide a wraparound defect}
On a four-cycle, take proposed differences \((2,-1,3,-3)\) and anchor zero. Integrate the first three edges. Which equation fails, and what changes when only the final edge is allowed to move?

\textbf{Solution.} The first three edges give values \((0,2,1,4)\). Their required last difference is \(x_0-x_3=-4\), not \(-3\). The total prescribed sum is one, so the original loop has no realization. Keeping the first three edges forces the unique correction \((2,-1,3,-4)\), whose sum is zero. This is a protected repair of the edge data. It is not a realization of the original data without modification.

\subsection*{6. Residual decay can hide unprotected drift}
In maximum-norm $\R^3$, use $p(s,u,r)=s$ and $d(s,u,r)=r$. Compare the paths $(1,0,(n+1)^{-1})$ and $(1,H_n,(n+1)^{-1})$. Which hypotheses distinguish them?

\textbf{Solution.} Both preserve $p=1$ and have the same residual ratios $(n+1)/(n+2)$. The first moves by $1/((n+1)(n+2))$ and converges to $(1,0,0)$. The second moves by $1/(n+1)$ and has an unbounded second coordinate. In the variable-rate theorem, the first has $C_n=1/(n+2)$ and a summable weighted budget. The second has $C_n=1$, making that budget harmonic and divergent. The difference is movement, not residual decay or semantic preservation.

\subsection*{7. How a finite observation changes the apparent rate}
Let $\eta_n=1/(n+1)$. Find its profile after sampling at $\phi(k)=2^k-1$. Does the new profile prove exponential convergence in the original index?

\textbf{Solution.} The sampled value is $\eta_{\phi(k)}=2^{-k}$. Recovering the original index gives $k=\log_2(n+1)$ on the sampled stages, so the value is still $1/(n+1)$. The schedule preserves vanishing but changes the meaning of the horizontal axis. No faster original process has been proved. KS75 is a convergence equivalence under a monotone cofinal schedule, not a rate-invariance theorem.

\subsection*{8. What does a zero floor fail to certify?}
For the diagonal model of Section~\ref{sec:escape}, can an exact protected answer exist in $\ell^2\oplus\R$? Does a coherent family of exact finite prefixes exist? Explain why the answers differ without contradicting finite-fiber existence.

\textbf{Solution.} Exactness with protected value one forces every sequence coordinate to equal minus one. This is not square summable, so there is no protected Hilbert answer. Each finite prefix nevertheless has the unique solution $(-1,\ldots,-1;1)$, and these solutions are coherent under truncation. Their inverse limit contains the all-minus-one sequence in the full product. The finite-fiber theorem supplies that coherent point. It does not prove that the point is in the image of the square-summable semantic space. The correction image is dense but does not contain the raw residual, so its floor is zero and its algebraic class is nonzero.

\subsection*{9. A phase witness without enumerating every gauge}
For $c(x,y)=x_2y_1$ on $\mathbf F_2^2$, show that no scalar coboundary equals $c$. Then compute the square of the lifted label $(0,e_1+e_2)$ in the eight-element extension.

\textbf{Solution.} Scalar coboundaries $a(x)+a(y)-a(x+y)$ are symmetric, but $c(e_1,e_2)=0$ and $c(e_2,e_1)=1$. Thus no scalar coboundary equals $c$. For $x=e_1+e_2$, $c(x,x)=1$, so $(0,x)^2=(1,0)$ and its fourth power is the identity. The corresponding matrix is $XZ$, whose square is $-I$. The additional phase bit records precisely the sign lost by projecting back to the four-label group.

\subsection*{A checklist for applying the theory}
Before selecting a theorem, state the raw and exact objects, the coefficient structure, the maps that relate views, and the data that must remain fixed. Then ask which conclusion is needed: existence, a description of all answers, an identification-preserving comparison, a stable distance estimate, or an executable finite certificate. Finally record the hypotheses that make that conclusion available.

The examples above fail only when one of those distinctions is suppressed. The missing information is not supplied by giving the same construction a more general name.

\clearpage
\chap{A hypothesis atlas}\label{app:hypotheses}
A longer monograph is useful only if the reader can locate the theorem that answers the actual question. This atlas groups the main obligations by the desired conclusion. The references lead to the statements and proofs, not just to examples that resemble the intended application.

\subsection{Existence and identification}
\begin{longtable}{@{}p{.24\linewidth}p{.48\linewidth}p{.21\linewidth}@{}}
\toprule Desired conclusion & Data and hypotheses to establish & Main entry\\\midrule\endhead
A coherent family & A view system with its actual restrictions and all compatibility equations. Nonemptiness is separate. & KS01--KS03\\[4pt]
A prescribed extension & The declared partial section and an actual vanishing relative obstruction, or matching data on a transition cover. & KS04, KS09, KS14\\[4pt]
A unique extension & Existence and vanishing of the invisible coherent ambiguity, not existence alone. & KS10, KS32\\[4pt]
An infinite point from finite views & Nonempty finite fibers, or another proved lifting/compactness hypothesis. A prescribed value must lie in the persistent image. & KS54, KS58\\[4pt]
Descent of objects and witnesses & The actual groupoid data, coherent overlap isomorphisms, and a composition-covering family on the specified site. & KS34--KS36\\[4pt]
Invariance of a boundary problem & An identified boundary and the specified comparison equivalences. Literal freezing of every representative is a different construction. & KS38--KS39, KS68\\[4pt]
The same limit from two paths & Convergence of both paths in a metric space and a pairwise distance budget tending to zero. & KS74\\
\bottomrule
\end{longtable}

\subsection{Approximation and numerical meaning}
\begin{longtable}{@{}p{.24\linewidth}p{.48\linewidth}p{.21\linewidth}@{}}
\toprule Desired conclusion & Data and hypotheses to establish & Main entry\\\midrule\endhead
An exact limit of approximations & Completeness, a nonnegative lower semicontinuous residual, vanishing residuals, and summable candidate movement. & KS71--KS72\\[4pt]
A protected exact limit & In addition, continuity of the protected map and preservation along the candidate path. & KS73\\[4pt]
Exact cancellation from zero floor & Closedness of the actual admissible residual image. Without it only closure membership follows. & KS76--KS77\\[4pt]
A stable residual-to-distance bound & The required closed range or bounded inverse in the declared normed model. In a protected problem, use the restriction to $\ker p$. & KS43, KS56\\[4pt]
A controlled moving exact space & Two residual certificates, or a proved split-surjective regime with the strict smallness condition. Norm closeness alone does not preserve rank. & KS48--KS53\\[4pt]
An independent object distance & An arrow category and an extended nonnegative cost with zero identity cost and composition subadditivity. Separation is another question. & KS78--KS79\\
\bottomrule
\end{longtable}

\subsection{Three questions before using a finite check}
First determine whether the tested object is the actual object in the theorem. A scalar simulation can illustrate an operator-valued statement without checking every dependent fiber. Next determine whether the test checks a universal finite domain or a selected finite range. Finally identify any limit that remains outside the computation.

The new phase extension is exhaustively checked on its eight elements. The diagonal example is different: rational finite sums verify its displayed enclosures, while the proof of nonexistence uses the infinite square-summability condition. Calling both results simply ``tested'' would conceal the most important distinction.

The companion evidence guide and \path{data/deepening_models.json} retain these scopes. The general source declarations are checked in the accepted formal baseline. Printed examples and new conceptual connections do not acquire additional Lean coverage merely because they occur next to a source identifier.

\chap{Notation and source concordance}\label{app:concordance}
\subsection{Notation}
\begin{longtable}{@{}p{.25\linewidth}p{.7\linewidth}@{}}
\toprule
Symbol & Meaning\\\midrule\endhead
\(\V,F_v,f^*\) & A declared view category, its fibers, and contravariant restrictions.\\
\(\Gamma(F)\) & The strict set of coherent global families.\\
\(\Gamma_U(F),\Ext_U(s)\) & Partial sections on a full patch and their global extension fiber.\\
\(\delta,\partial,d,D\) & Respectively the view differential, tower differential, a model residual, and a Hilbert defect operator. Each target is declared where used.\\
\(p,p_U,p_L\) & A protected observable, restriction to prescribed views, and a splitting dependent innovation observable.\\
\(\mathfrak o_U,H_U\) & The existence obstruction and the group of invisible coherent changes. Additive and nonabelian conventions are stated separately.\\
\(C_j,J_j,Q_N,S_N\) & Coarsener, chosen lift, restored coarse remainder, and retained finite reconstruction.\\
\(c_f,a_v\) & Coherent nonabelian transition labels and gauge values.\\
\(\mathfrak K_G,J_{\rm comp}\) & The cocycle patch groupoid and its composition-support topology.\\
\(\mathcal E_b(R)\) & Realizations equipped with an actual identification of their boundary with \(b\).\\
\(T_i,K_{ij},R_p\) & Transport to the root, the corresponding root loop operator, and reconstruction from the root.\\
\(\mathcal H,q,P\) & Root residual, an orthogonal kernel projector, and the explicitly defined model repair. A graph repair need not be orthogonal in the global maximum norm.\\
\(\pi_{W,U},H_U\) & The map forgetting part of the prescription and the corresponding invisible coherent ambiguity.\\
\(\phi,\operatorname{EventualHom}\) & A monotone source schedule and transport classes modulo eventual equality.\\
\(\chi_A,M=SA-BC\) & Split exact/residual coordinates and the mismatch between chosen repair lifts.\\
\(\omega,Q_n,C_n,b_n\) & A residual, its prefix-product rate, movement coefficient, and summable step budget. Here $Q_n$ is not the dyadic remainder operator $Q_N$.\\
\(\nu_B,d_W,d_W^{\mathrm{sym}}\) & Residual floor, extended directed transport distance, and its symmetric maximum.\\
\(c(x,y),\kappa(x,y)\) & The two-bit pair phase and its antisymmetric scalar-gauge invariant.\\
\(\simeq,\simeq_L\) & A specified equivalence, and a continuous linear equivalence. Neither notation alone asserts an isometry.\\
\bottomrule
\end{longtable}

\subsection{How to read the correspondence}
The identifiers below are stable editorial identifiers, not counts of new discoveries. A row can combine several source declarations into one exposition. The named declaration, its line range, source hash, and correspondence note are recorded in \code{data/claim_ledger.json} and in the companion Lean guide. Existence of a matching source declaration is not a mechanical proof that every word of an edited statement matches its type. The displayed hypotheses and the printed proof remain part of the review.

\textbf{F} denotes the recovered foundation tree at version 0.20.0. \textbf{M} denotes the focused moving-space tree at version 0.32.0. \textbf{S} denotes the separately recovered Hilbert holonomy source. \textbf{I} denotes a source in the canonical integrated tree, whether inherited or added during the version 1.2.0 formalization. F, M, and S retain provenance, while byte-identical copies of their referenced sources were compiled together in the accepted version 1.2.0 replay of \path{formal/KomorebiSpace/}. All \EditionClaims\ principal items have source correspondences in that checked closure. This does not assert that every illustrative computation or explanatory sentence is independently formalized.

\input{claim_table}

\par\noindent\begin{minipage}{\linewidth}\subsection{Verification is a separate record}
The foundation archive reports 122 production proof modules, 1,299 named theorem declarations, and 696 indexed constructions. Its recorded strategy includes one complete cumulative source build, followed by a fresh archive replay that reuses a checked 117-module baseline and compiles the five additions and root again. It must not be described as two full clean cumulative builds.

The focused moving-space archive reports 11 production modules and a root, 151 named theorem declarations, and 47 indexed constructions. Its historical replay record reports two complete clean builds of that focused source closure. The archive SHA-256 and the retained source hashes are checked by this publication's integrity script. The focused tree does not replace the whole historical Komorebi Space repository.

The pinned environment in the records is Lean 4.34.0-rc2 with mathlib commit\par\smallskip\noindent\path{aeb72f1e9981b6d35b426b16dca7759be3e7ba06}.\par\smallskip The permitted standard axioms are \code{propext}, \code{Classical.choice}, and \code{Quot.sound}. The compiler, kernel, and cached dependencies remain trusted. For version 1.2.0 the portable runtime was restored and its version and dependency pins checked, but the compiler and mathlib were not rebuilt independently. This expanded publication retains that verified baseline without executing a new compiler replay.

The retained version 1.2.0 formal record covers \EditionFormalModules\ production modules and one root. Clean source replays and an imported axiom audit are separate from archive integrity, exact finite computations, and PDF preflight. The five modules added in that formalization contain \EditionNewTheorems\ explicit named theorems and lemmas, including supporting results. Counts are not claims of mathematical novelty. The complete source hashes, command logs, and baseline scope are in \path{audit/formal/release.json}.

\end{minipage}

\subsection*{Editorial source policy}
Earlier source documents contain historical scope statements and counts. Later checked developments can supersede an earlier statement that a topic is a target. In particular the old outline's stack target is not the status of the later standard stack theorem. This monograph uses the actual relevant declarations rather than treating the earliest outline as a current completion ledger.



\subsection*{Additional source connections in this edition}
The earlier expansion developed KS59--KS70. KS71--KS80 add exactification, floors, transport costs, and the phase obstruction. Their declarations are listed above. Separately scoped worked models appear in \path{data/deepening_models.json}. No new compiler execution is claimed.

\input{chapters/08_references}
\end{document}
